% UBT-AI-PROVENANCE-BEGIN
% schema: ubt-ai-provenance/v1
% tier: C_working
% ai_assistance: disclosed
% human_review: risk-based
% editorial_responsibility: Ing. David Jaroš
% policy: ../../../AI_PROVENANCE.md
% notice: Working material; exhaustive human review is not claimed.
% UBT-AI-PROVENANCE-END

\chapter{Spekulativní rozšíření (číst odděleně)}

\section{Hranice tohoto dodatku}
Tento dodatek je mapou výzkumných směrů, nikoli zdrojem kanonických axiomů.
Žádný z následujících bodů nesmí být použit k uzavření mezery v hlavní
teorii, dokud nejsou jeho předpoklady, rovnice a redukční zobrazení nezávisle
odvozeny.

\section{Otázky komplexního času a fyzické fáze}
Kanonickou souřadnicí je $\tau=t+i\psi$. Možná kompaktnost, tepelná periodicita
nebo dynamická stabilizace směru $\psi$ zůstávají výzkumnými otázkami.
Preferovaný řez ani průměrování přes vlákno se nesmí vkládat do lokální
definice metriky jen proto, aby se zjednodušilo pozorování.

\section{Topologické a solitonové sektory}
Hopfionová, vinutá a defektní řešení lze studovat jako kandidátní sektory
zadané akce. Samotná topologie neurčuje hmotnosti částic, vazby ani stabilitu.
Každý návrh potřebuje výsledek o existenci, podmínku konečné energie a
analýzu fluktuací.

\section{Modely související s vědomím}
Modely, které spojují struktury fázového prostoru nebo stochastické struktury
s vědomím, patří do spekulativní výzkumné vrstvy. Nejsou důsledkem
kanonických axiomů UBT a v současnosti nemají zavedený empirický most.

\section{Kontrolní seznam pro povýšení}
Spekulativní konstrukce může vstoupit do aktivní fyzikální linie teprve
poté, co uvede falzifikovatelnou pozorovatelnou veličinu, odvodí své proměnné
z architektury jediného pole, zachová zamčenou geometrii a projde
repozitářovými kontrolami verifikace a dvojjazyčné revize.
